\section{Conclusión Profesional (Pregunta 4)}

\textit{Si fueras el ingeniero consultor de una empresa logística, ¿en qué
situación recomendarías invertir tiempo computacional en NSGA-II para mostrar
el frente completo, y en qué situación usarías un enfoque directo a priori?}

\subsection{El criterio de decisión: ¿conoces tus pesos?}

La pregunta no es ``¿qué algoritmo es mejor?'', sino ``¿qué tan seguro está el
negocio de sus preferencias \emph{antes} de optimizar?''. Mi experiencia en
banca (Santander US) lo ilustra con dos analogías directas:

\paragraph{Caso A --- Decisión estratégica única (recomiendo NSGA-II).}
Es análogo a definir la estrategia de \textbf{reabastecimiento de cajeros
(ATMs)}: el costo logístico del traslado de efectivo contra el SLA de
disponibilidad ante el cliente. Cuando la decisión es de alto impacto, se toma
una sola vez, involucra a múltiples \emph{stakeholders} (Operaciones, Riesgo,
Experiencia del Cliente) y \textbf{nadie tiene pesos claros consensuados},
vale la pena invertir el cómputo de NSGA-II. Mostrar el frente de Pareto
completo convierte una discusión política (``¿cuánto vale un minuto de
disponibilidad?'') en una negociación informada sobre compromisos concretos y
visibles. El frente es una herramienta de \emph{alineación organizacional},
no sólo de optimización.

\paragraph{Caso B --- Decisión operativa recurrente (recomiendo A Priori).}
Es análogo a la \textbf{asignación automática de líneas de crédito} o al
ruteo diario de reabastecimiento: decisiones que se ejecutan miles de veces,
con pesos ya \emph{aprendidos y estables} (el negocio sabe que prioriza, por
ejemplo, 70\,\% costo / 30\,\% tiempo). Aquí desplegar NSGA-II en cada corrida
sería un desperdicio: el enfoque a priori entrega en una fracción del tiempo
una solución que ---como demostró la Sección 4--- es prácticamente idéntica a
la que se elegiría sobre el frente completo, gracias a la convexidad del
problema.

\subsection{Síntesis de hallazgos}

\begin{itemize}[noitemsep]
  \item \textbf{Suma ponderada:} a priori y a posteriori \emph{convergen}
        (diferencia de sólo $\num{\DiffPondCosto}$ en costo) cuando el frente
        es convexo. El cómputo extra de NSGA-II no cambia la decisión.
  \item \textbf{Lexicográfica:} la prioridad rígida sobre $f_2$ disparó el
        costo a $f_1 = \num{\AprioriLexCosto}$ (el peor del frente). Útil sólo
        ante objetivos innegociables.
  \item \textbf{Valor del frente:} su retorno no es numérico sino de
        \emph{transparencia}: expone el menú completo de compromisos.
\end{itemize}

\subsection{Recomendación final}

\begin{center}
\fbox{\begin{minipage}{0.92\linewidth}
\textbf{Regla práctica.} Usa \textbf{A Posteriori (NSGA-II)} cuando la decisión
sea estratégica, única, con \emph{stakeholders} múltiples y sin pesos
consensuados: el frente alinea a la organización. Usa \textbf{A Priori (GA)}
cuando la decisión sea operativa, recurrente y con preferencias estables ya
conocidas: es más barato y, sobre frentes convexos, igual de acertado.
\end{minipage}}
\end{center}

En definitiva, el momento de expresar las preferencias debe alinearse con la
\textbf{naturaleza de la decisión}, no con la sofisticación del algoritmo
disponible. Un buen consultor elige la herramienta proporcional al problema.
